\documentclass[../cours_2de.tex]{subfiles}
\begin{document}
\setcounter{chapter}{15}
\chapter{Variations, extrémums et optimisation}
\label{sec:ch16}
\begin{objectif}
Lire et construire un tableau de variations ; distinguer observation et preuve ; résoudre des problèmes d'optimisation sans dérivées.
\end{objectif}
\section{Les définitions précises}
\begin{definition}
$f$ est croissante sur un intervalle $I$ si, pour tous $a,b\in I$ tels que $a<b$, on a $f(a)\leq f(b)$. Elle est strictement croissante si l'inégalité des images est stricte.
Elle est décroissante si $f(a)\geq f(b)$ dans la même situation. Une fonction monotone est croissante ou décroissante sur l'intervalle considéré.
\end{definition}
Pour établir une croissance, il faut comparer deux entrées quelconques, pas seulement deux points du graphique.
Un tableau de variations résume les intervalles où la monotonie est connue, les valeurs aux bornes lorsqu'elles sont définies et les extrémums atteints. Une flèche n'indique pas une pente constante.
\begin{exemple}
Pour $f(x)=(x-2)^2+1$ sur $[-1;4]$ :
\[
\begin{array}{c|ccccc}
x&-1&&2&&4\\\hline
f(x)&10&\searrow&1&\nearrow&5
\end{array}
\]
Le minimum vaut $1$, atteint en $2$ ; le maximum sur cet intervalle vaut $10$, atteint en $-1$. Pour deux entrées situées de part et d'autre de $2$, le tableau seul ne fixe pas l'ordre des images.
\end{exemple}
\section{Un extremum est une valeur atteinte}
\begin{definition}
$f$ admet un maximum $M$ sur $D$ s'il existe $a\in D$ tel que $f(a)=M$ et si, pour tout $x\in D$, $f(x)\leq M$.
Un minimum $m$ est défini avec $f(x)\geq m$ et une valeur où l'égalité est atteinte.
\end{definition}
\begin{exemple}
Sur $]0;1[$, $f(x)=x$ reste entre $0$ et $1$ mais n'atteint aucun des deux : elle n'a ni minimum ni maximum. Sur $[0;1]$, elle atteint les deux.
Montrer seulement $f(x)\leq10$ ne prouve pas que $10$ est son maximum.
\end{exemple}
Un tableau de valeurs ou une courbe tracée par un logiciel peut suggérer des variations ; il peut manquer un pic étroit entre deux points.
De même, une fonction croissante peut avoir une discontinuité : la monotonie seule ne garantit pas qu'elle prenne toutes les valeurs intermédiaires. Nous n'invoquons cette dernière propriété que pour les fonctions continues usuelles, avec son statut d'outil admis.
\section{Prouver sans dériver}
\begin{methode}
Comparer $f(b)-f(a)$ pour $a<b$ ; reconnaître une fonction de référence après un changement de variable ; ou écrire une quantité sous forme d'un carré positif ou nul.
Pour un extremum, établir une borne générale puis montrer qu'elle est atteinte.
\end{methode}
\begin{exemple}
$g(x)=-3(x-1)^2+7\leq7$, avec égalité exactement pour $x=1$.
Sur $]-\infty;1]$, le carré $(x-1)^2$ décroît ; multiplier par $-3$ renverse l'ordre, donc $g$ croît. Sur $[1;+\infty[$, elle décroît.
\end{exemple}
\section{Optimiser une aire}
\begin{exemple}[Rectangle de périmètre fixé]
Pour un périmètre de $20$ mètres, les côtés sont $x$ et $10-x$, avec $0<x<10$.
L'aire est
$A(x)=x(10-x)=25-(x-5)^2\leq25$.
L'égalité est atteinte pour $x=5$, valeur admissible. L'aire maximale vaut donc $25\unite{m^2}$ et correspond à un carré de côté $5$ mètres.
\end{exemple}
On peut aussi utiliser l'inégalité arithmético-géométrique :
$\sqrt{x(10-x)}\leq(x+10-x)/2=5$, puis élever au carré les quantités non négatives.
Cela fournit une autre preuve du même résultat.
\begin{attention}
Si le domaine était $0<x<4$, la valeur $x=5$ ne serait pas admissible. Sur cet intervalle, l'aire croît mais ne possède pas de maximum : la borne $24$ est approchée quand $x$ se rapproche de $4$, sans être atteinte.
\end{attention}
\section{Approximer : balayage et dichotomie}
Un balayage évalue $f$ sur une grille et retient la meilleure valeur. Il calcule un maximum \emph{sur la grille}, pas automatiquement sur le domaine continu.
\begin{exemple}
Pour $F(x)=2-|x^2-2|$ sur $[0;2]$, $F(x)=x^2$ jusqu'à $\sqrt2$, puis $F(x)=4-x^2$.
Le tableau de variations montre un maximum $2$ atteint en $\sqrt2$.
Pour en approcher l'abscisse, maintenir un encadrement $a\leq\sqrt2\leq b$, tester $m=(a+b)/2$, puis remplacer $a$ par $m$ si $m^2<2$, sinon $b$ par $m$.
\end{exemple}
\begin{demonstration}[Garantie de la dichotomie]
La croissance du carré sur les positifs assure que l'intervalle conservé contient toujours $\sqrt2$.
Sa longueur est divisée par $2$ à chaque étape : depuis $[1;2]$, elle vaut $2^{-n}$ après $n$ étapes. Dix étapes donnent une amplitude inférieure à $10^{-3}$.
\end{demonstration}
Ce procédé localise un extremum grâce à une propriété spécifique démontrée ; ce n'est pas un algorithme universel de maximisation d'une fonction quelconque.
\begin{finchapitre}
\section{Approcher une longueur de courbe}
Pour $x_0<\cdots<x_n$, on relie les points $(x_i,f(x_i))$ par des segments. La longueur polygonale est
\[
L_n=\sum_{i=0}^{n-1}\sqrt{(x_{i+1}-x_i)^2+(f(x_{i+1})-f(x_i))^2}.
\]
La formule vient de la distance entre deux points et de l'addition des longueurs. Pour une courbe régulière comme $y=x^2$, raffiner la subdivision donne une approximation de sa longueur ; l'existence et la convergence vers cette longueur relèvent d'une étude ultérieure.
\utilite{Optimiser signifie chercher la meilleure valeur admissible et la justifier. Séparer preuve, modèle et approximation sera essentiel en analyse et en calcul scientifique.}
\begin{retenir}
Un extremum doit être atteint. Toujours donner le domaine, une borne et un cas d'égalité. Un résultat numérique ne suffit pas à établir une propriété globale.
\end{retenir}
\end{finchapitre}
\end{document}
